\section{Introduction}

In Griggs v. Duke Power Co.~\cite{griggs1971}, the US Supreme Court ruled a business hiring decision illegal if it resulted in disparate impact by race even if the decision was not explicitly determined based on race.  The Duke Power Co. was forced to stop using intelligence test scores and high school diplomas, qualifications largely correlated with race, to make hiring decisions. The Griggs decision gave birth to the legal doctrine of \emph{disparate impact}, which today is the predominant legal theory used to determine \emph{unintended} discrimination in the U.S. Note that disparate \emph{impact} is different from disparate \emph{treatment}, which refers to \emph{intended} or direct discrimination.  Ricci v. DeStefano \cite{Ricci09} examined the relationship between the two notions, and disparate impact remains a topic of legal interest.

\looseness-1 Today, algorithms are being used to make decisions both large and small in almost all aspects of our lives, whether they involve mundane tasks like recommendations for buying goods, predictions of credit rating prior to approving a housing loan, or even life-altering decisions like sentencing guidelines after conviction\cite{Hodson15NewScientist}. How do we know if these algorithms are biased, involve illegal discrimination, or are unfair? 

%  As creators of computational tools that make unsupervised decisions, we believe it behooves us, both legally and ethically, to study disparate impact under a computational lens.
% %
% Legally, we want to understand whether and how present data mining practices meet the criteria that characterize disparate impact, and whether the criteria are robust with respect to variations of practice.
% %
% Ethically, we ask whether data mining can achieve unequitable outcomes, how to interpret these ideas computationally, and whether current tools can be improved.
%
These concerns have generated calls, by governments and NGOs alike, for research into these issues~\cite{podesta14bigdata, civilRightsPrinciples}.
In this paper, we introduce and address two such problems with the goals of quantifying and then removing disparate impact.

%\vspace*{0.1in}\mypara{Disparate Impact.}
While the Supreme Court has resisted a ``rigid mathematical formula'' defining disparate impact~\cite{watson1988}, we will adopt a generalization of the 80 percent rule advocated by the US Equal Employment Opportunity Commission (EEOC)~\cite{eeoc1979}. 
We note that disparate impact itself is not illegal; in hiring decisions, business necessity arguments can be made to excuse disparate impact.%\footnote{For example, in New York City Transit Authority v. Beazer the transit authority successfully made the case that they should be allowed to bar drug users (including those using drugs to aid the process of recovery) from employment~\cite{Barocas14DisparateImpact}.}
\begin{defn}[{\small Disparate Impact (``80\% rule'')}]
Given data set $D = (X, Y, C)$, with \emph{protected attribute} $X$ (e.g., race, sex, religion, etc.), remaining attributes $Y$, and binary class to be predicted $C$ (e.g., ``will hire''), we will say that $D$ has \emph{disparate impact}
if 
\[ \frac{\Pr(C = YES| X = 0)}{\Pr(C = YES| X = 1)} \leq \tau = 0.8 \]
for positive outcome class $YES$ and majority protected attribute $1$ where $\Pr(C = c | X = x)$ denotes the conditional probability (evaluated over $D$) that the class outcome is $c \in C$ given protected attribute $x \in X$.\footnote{Note that under this definition disparate impact is determined based on the given data set and decision outcomes. Notably, it does not use a broader sample universe, and does not take into account statistical significance as has been advocated by some legal scholars~\cite{peresie09law}.}
\end{defn}

\looseness-1 The two problems we consider address identifying and removing disparate impact.  The \emph{disparate impact certification problem} is to guarantee that, given $D$, any classification algorithm aiming to predict some $C'$ (which is potentially different from the given $C$) from $Y$ would not have disparate impact.  By certifying any outcomes $C'$, and not the process by which they were reached, we follow legal precedent in making no judgment on the algorithm itself, and additionally ensure that potentially sensitive algorithms remain proprietary.  The \emph{disparate impact removal problem} is to take some data set $D$ and return a data set $\bar{D} = (X, \bar{Y}, C)$ that can be certified as not having disparate impact.  The goal is to change only the remaining attributes $Y$, leaving $C$ as in the original data set so that the ability to classify can be preserved as much as possible.

\subsection{Results}

We have four main contributions.

We first introduce these problems to the computer science community and develop its theoretical underpinnings.  The study of the EEOC's 80\% rule as a specific class of loss function does not appear to have received much attention in the literature.  We link this measure of disparate impact to the balanced error rate (\ber).
We show that any decision exhibiting disparate impact can be converted into one where the protected attribute leaks, i.e. can be predicted with low \ber.

Second, this theoretical result gives us a procedure for certifying the impossibility of disparate impact on a data set.  This procedure involves a particular regression algorithm which minimizes \ber. We connect \ber to disparate impact in a variety of settings (point and interval estimates, and  distributions). We discuss these two contributions in Sections \ref{sec:note-meas-disp} and \ref{sec:main-idea}.

In Section \ref{sec:repairing-data}, we show how to transform the input dataset so that predictability of the protected attribute is impossible. We show that this transformation still preserves much of the signal in the unprotected attributes and has nice properties in terms of closeness to the original data distribution.

Finally, we present a detailed empirical study in Section \ref{sec:experiments}.  We show that our algorithm certifying lack of disparate impact on a data set is effective, such that with the three classifiers we used certified data sets don't show disparate impact.  We demonstrate the fairness / utility tradeoff for our partial repair procedures.  Comparing to related work, we find that for any desired fairness value we can achieve a higher accuracy than other fairness procedures.  This is likely due to our emphasis on changing the data to achieve fairness, thus allowing any strong classifier to be used for prediction.

Our procedure for detecting disparate impact goes through an actual classification algorithm. As we show in our experiments, a better classifier provides a more sensitive detector.
We believe this is notable. As algorithms get better at learning patterns, they become more able to introduce subtle biases into the decision-making process by finding subtle dependencies among features. But this very sophistication helps \emph{detect} such biases as well via our procedure! Thus, data mining can be used to verify the fairness of such algorithms as well.

\section{Related Work}
There is, of course, a long history of legal work on disparate impact.  There is also related work under the name \emph{statistical discrimination} in Economics.  We will not survey such work here.  Instead, we direct the reader to the survey of Romei and Ruggieri~\cite{Romei13Multidisciplinary} and to a discussion of the issues specific to data mining and disparate impact~\cite{Barocas14DisparateImpact}.  Here, we focus on data mining research relating to  combating discrimination.  This research can be broadly categorized in terms of methods that achieve fairness by modifying the classifiers and those that achieve fairness by modifying data. 

 Kamishima et al.~\cite{Kamishima12Fairness, Kamishima11Fairness} develop a regularizer for classifiers to penalize prejudicial outcomes and show that this can reduce indirect prejudice (their name for implicit discrimination like disparate impact) while still allowing for accurate classification.  They note that as prejudicial outcomes are decreased, the classification accuracy is also decreased.
Our work falls into the category of algorithms that change the input data.  Previous work has focused on changing the class values of the original data in such a way so that the total number of class changes is small~\cite{Kamiran09Classifying, Calders09Building}, while we will keep the class values the same for training purposes and change the data itself.  Calders et al.~\cite{Calders09Building} have also previously examined one method for changing the data in which different data items are given weights and the weights are adjusted to achieve fairness.  In this category of work, as well, there is worry that the change to the data will decrease the classification accuracy, and Calders et al. have formalized this as a fairness/utility tradeoff~\cite{Calders09Building}.  We additionally note that lower classification accuracy may actually be the desired result, if that classification accuracy was due to discriminatory decision making in the past.

An important related work is the approach of ``fairness through awareness'' of Dwork et al.~\cite{Dwork12Fairness} and Zemel et al.~\cite{icml2013_zemel13}. Dwork et al.~\cite{Dwork12Fairness} focus on the problem of \emph{individual} fairness; their approach posits the existence of a similarity measure between individual entities and seeks to find classifiers that ensure similar outcomes on individuals that are similar, via a Lipschitz condition. In the work of Zemel et al.~\cite{icml2013_zemel13}, this idea of protecting individual fairness is combined with a statistical group-based fairness criterion that is similar to the approach we take in this work. A key contribution of their work is that they learn a modified representation of the data in which fairness is ensured while attempting to preserve fidelity with the original classification task.  While this group fairness measure  is similar to ours in spirit, it does not match the legal definition we base our work on. Another paper that also (implicitly) defines fairness on an individual basis is the work by Thanh et al.~\cite{luong2011knn}. Their proposed repair mechanism changes class attributes of the data (rather than the data itself). 

Pedreschi, Ruggieri and Turini \cite{pedreschi2012study,pedreschi2009integrating} have examined the ``80\% rule'' that we study in this paper as part of a larger class of measures based on a classifier's confusion matrix. 

%\subsection{Legal Background}
%\label{sec:legal}
%Our understanding of the legal process is as follows \sfnote{cite Solon here}: disparate impact and the associated policy to be questioned are determined and challenged together, but are separate questions in the sense that a policy can be found to have disparate impact and still be determined to be legal.  Interestingly, because of the joint nature of such a legal challenge, disparate impact cases are not brought when the policy (algorithm) is unknown or overly complex.  These are exactly the cases which we are interested in addressing.  We believe that the complexity of an algorithm does not excuse it from causing disparate impact.
%One strength of the approach we present here is that complex or proprietary algorithms can now be repaired so that they no longer cause disparate impact, even without understanding or revealing the algorithm itself.  Our definition of the disparate impact removal problem ensures that such algorithms remain hidden and can change without introducing disparate impact as long as their input, $Y$, remains the same.  We hope that this will create clarity around the disparate impact of opaque algorithms, and serve as a mechanism by which they can be challenged as policy.
%
%We will use hiring as our main motivating example for Section \ref{sec:main-idea} and college admissions (affirmative action) for Section \ref{sec:repairing-data}, but these results generalize to many other data sets and domains in which data mining techniques have and will become prevalent and protected attributes should not be used to make decisions.  Interestingly, while we see a direct relationship between the finding of disparate impact in college admissions decisions and the repair of that impact by affirmative action, the complexity of the college admissions process (or admissions algorithm) means that it is hard to challenge the process itself in court.  Instead, affirmative action policies based on race have been challenged as having  \emph{disparate treatment}, or explicitly racially biased outcomes.  Ironically, this is the legal danger of attempting to repair disparate impact --- if your repair is explicitly based on a protected attribute, it may be found to be illegal (as was the case in Ricci v. DeStefano~\cite{Ricci09}).  While we will suggest a repair algorithm in Section \ref{sec:repairing-data} based on protected attributes, any policy implementing this algorithm would need to consider the specific choice of attribute with caution in the legal context of the application.  One such successful example is the Texas Top Ten Percent Rule~\cite{Texas97Top10Percent}, which admits the top ten percent of every high school class in Texas to the University of Texas.
%By using an admissions algorithm repair based on high school class, which is explicitly racially neutral but highly correlated with race, the Top Ten Percent Rule avoids disparate treatment while addressing disparate impact.  The disparate impact removal algorithm we introduce here is a general solution that can be implemented for any chosen attribute --- we will leave such policy decisions for others.


%%% Local Variables:
%%% mode: latex
%%% TeX-PDF-mode: t
%%% TeX-master: "bias"
%%% End:
